% ID: SEL-D0017-BA-E01
% ESTADO: cotejada. Segunda pasada visual completa y borrador compilado revisados.
% Fuente: Selectividad/matematicas_aplicadas/matematicas_aplicadas_P_E.pdf
% SHA-256: bc88bfeb0f184eb0c806c4efe04e84c24f3e5fcef9caaa12e94fd48d9f67a389
% Página: 1; bloque A; ejercicio original 1.

\textbf{(2.5 puntos)} Una confitería elabora dos tipos de tartas, unas de
chocolate y otras de merengue y chocolate. Para ello dispone de 100 kg de
bizcocho, 80 kg de crema de chocolate y 46 kg de merengue. Para elaborar
una tarta de chocolate, se requieren 1 kg de bizcocho y 2 kg de crema de
chocolate y para la tarta de chocolate y merengue se requieren 2 kg de
bizcocho, 1 kg de crema de chocolate y 1 kg de merengue. Por cada tarta
de chocolate se obtiene un beneficio de 10 euros y de 12 euros por cada
una de merengue y chocolate. Suponiendo que se vende todo lo que se elabora,
\textquestiondown{}cuántas tartas de cada tipo debe preparar para obtener
un beneficio máximo? \textquestiondown{}Cuál es dicho beneficio?
